% output=pdftex

\starttext

\usemodule[abr-02]

\environment corrlogo % also defines color

\defineoverlay [page] [\reuseMPgraphic{page}]

\definefilesynonym [correspondence] [x-corres.rng]

\startsetups rng:contact
  \showXMLwrd[ref]
  \showXMLlin[optional]
\stopsetups

\startsetups rng:fields
  \showXMLtxt[contact,contactgroup]
  \showXMLlin[initials,formalname,informalname,title,prefix,p,member]
\stopsetups

\setuplayout
  [height=middle,
   width=middle,
   header=2cm,
   footer=2cm,
   topspace=.05\paperheight,
   bottomspace=.05\paperheight,
   backspace=.1\paperwidth,
   cutspace=.1\paperwidth]

\setuptyping
  [option=color]

\usetypescriptfile[type-fsf]

\definetypeface [mainface] [ss] [sans]  [opus]     [default] [encoding=texnansi]
\definetypeface [mainface] [rm] [serif] [palatino] [default] [encoding=texnansi]
\definetypeface [mainface] [mm] [math]  [palatino] [default]
\definetypeface [mainface] [tt] [mono]  [modern]   [default] [rscale=1.1]

\setupbodyfont [mainface,ss]

\definefontsynonym [MainNormal] [Sans]
\definefontsynonym [MainBold]   [SansBold]

\setuphead[chapter][style=\bfc,color=DarkGray]
\setuphead[section][style=\bfa,color=DarkGray]

\setupwhitespace[big]

\setupbackgrounds [page] [background=page]

\startstandardmakeup[color=LightGray]

  \vfill

  \WidthSpanningText
    {correspondence\setstrut\strut}
    {\textwidth}
    {SansBoldCaps}

  \midaligned{\WidthSpanningText{Hans Hagen}{.5\textwidth}{MainBold}}

\stopstandardmakeup

\setupbackgrounds [page] [background=]

\setuphead[chapter][page=,before={\blank[6*medium]},after={\blank[2*medium]}]

\chapter{Introduction}

This manual describes how to handle correspondence with
\CONTEXT. The core components are a simple \XML\ database
(the successor of a \TEX\ based format), a few styles, and
a graphical user interface. It is a nice example of combined
technologies.

\chapter{The Database}

One can argue about the structure of this database, but
since it suits our purpose, we will not spend much words
about how it ended up this way. You can happily use your own
database and let it produce this kind of \XML\ code, which
is then happily accepted by \CONTEXT.

As said, the structure is quite similar to the \TEX\
counterpart that we have been using for years to typeset
our letters, envelopes, mailings and alike. This is
typically the kind of code that remains unchanged for
years, but the advance of \XML\ and evolution of the
\EXAMPLE\ framework triggered a rewrite.

The root element is defined as follows (we use \RELAXNG\ as
definition format).

\showRNGcomponent[correspondence][contacts]

Each contact is a collection of data fields which may come
in any order.

\start
  \setups[rng:contact]
  \showRNGcomponent[correspondence][contact]
\stop

A contact group is a collection of references to contacts:

\showRNGcomponent[correspondence][contactgroup]

A contact file is just a reference to a to be included file:

\showRNGcomponent[correspondence][contactfile]

The fields that make up a contact are normally just text.
The address needs to be preformatted with \type {<p>}
elements (more details can be found in the file \type
{x-corres.rng}).

A simple database will look like this:

\startbuffer
<contacts>

<contact label='hagen.j'>
  <initials>J.</initials>
  <formalname>Hagen</formalname>
  <informalname>Hans</informalname>
  <title>Drs.</title>
  <prefix>Heer</prefix>
  <address>
    <p>J. Hagen</p>
    <p>PRAGMA ADE</p>
    <p>Ridderstraat 27</p>
    <p>8061GH Hasselt NL</p>
  </address>
</contact>

<contact label='otten.a.f'>
  <initials>J.</initials>
  <formalname>Otten</formalname>
  <informalname>Ton</informalname>
  <title>Drs.</title>
  <prefix>Heer</prefix>
  <address>
    <p>A.F. Otten</p>
    <p>PRAGMA ADE</p>
    <p>Ridderstraat 27</p>
    <p>8061GH Hasselt NL</p>
  </address>
</contact>

<contact label='marle.k.van'>
  <initials>K. van</initials>
  <formalname>Marle</formalname>
  <informalname>Kees</informalname>
  <title>Ir.</title>
  <prefix>Heer van</prefix>
  <address>
    <p>K. van Marle</p>
    <p>PRAGMA POD</p>
    <p>Ridderstraat 27</p>
    <p>8061GH Hasselt NL</p>
  </address>
</contact>

<contactgroup label='pragma-ade'>
  <member>hagen.j</member>
  <member>otten.a.f</member>
</contactgroup>

<contactgroup label='pragma-pod'>
  <member>hagen.j</member>
  <member>otten.a.f</member>
  <member>marle.k.van</member>
</contactgroup>

</contacts>
\stopbuffer

\start
  \setups[rng:fields]
  \showXMLbuffer
\stop

The labels plays an important role later on, when we will
filter data from the database, so they should be unique and
consistent.

\chapter{The Interface}

You can use the database from within \CONTEXT\ after you've
loaded the file \type {x-corres.tex}:

\starttyping
\usemodule[corres]
\stoptyping

The macro that deals with processing is:

\starttyping
\XMLprocesscontacts[filename]
\stoptyping

The argument is optional, in which case the selection needs
to be set up with the following command:

\starttyping
\setvariables
  [contacts]
  [selection=,
   file=]
\stoptyping

This command is also used for specifying the data file.

The fields of a contact are stored and can be used in the
typesetting process. This process is hooked into the \XML\
engine and style by the following setup (here we show the
default definition):

\starttyping
\startsetups[contact:handle]

  \XMLflush{address}

\stopsetups
\stoptyping

And so we have arrived at the main process, which may look
like something:

\starttyping
\setvariables [contacts] [selection=pragma-ade,file=pragma.xml]

\startsetups [contact:handle]
  \XMLflush{initials} \XMLflush{formalname} \endgraf
\stopsetups

\XMLprocesscontacts
\stoptyping

A selection can be a comma separated list of labels and|/|or
groups. In a first pass, the list is resolved into a list of
valid contacts, while the second pass typesets the result
based on the handle.

\chapter{The Style}

In the previous chapters we presented the \XML\ database
and its interface, here will will spend some words on
typesetting a (series of) letter(s). The style is quite
independent of the database but uses the same fields and is
called \type {m-letter.tex}.

It does not make sense to discuss all the details of the
style because you need to look into the file anyway, so we
stick to discussing the methodology used. There is an
\EXAMPLE\ interface that permits you to play with the
settings.

The layout setup is rather simple and uses standard
\CONTEXT\ setup commands. The individual components that
make up the letter are defined in so called setups, like:

\starttyping
\startsetups [letter:place:address]
  ...
\stopsetups
\stoptyping

We use the layering mechanism to collect and place the
components. Inside such a setup, variables are used to
control the behaviour. For instance, the location of the
address on the page grid is determined by:

\starttyping
\setvariables
  [letter:address]
  [line=3]
\stoptyping

The address itself again a setup:

\starttyping
\setsetups[letter:address]
  ...
\stopsetups
\stoptyping

and although a definition like:

\starttyping
\startsetups[letter:address]
  A Fancy Name \endgraf
  A Nice Address \endgraf
  The Place To Go
\stopsetups
\stoptyping

is okay, it makes more sense to use variables again, like
in:

\starttyping
\startsetups[letter:address]
  \def\\{\endgraf}\getvariable{letter:data}{address}
\stopsetups
\stoptyping

Of course you need to define the adress:

\starttyping
\setvariables
  [letter:data]
  [address={Name\\Address\\City}]
\stoptyping

This chaining of settings may look overly complicated, but
the advantage is that by using setups we can easily
overload components.

For the same reason, we have isolated the labels for
instance:

\starttyping
\setuplabeltext
  [en]
  [letter:opening:formal=Dear,
   letter:opening:informal=Dear,
   letter:opening:unknown=LS]
\stoptyping

Labels are defined per language, so that the style
automatically adapts itself.

Now, when you load the style, and optionally overload setups,
and definitely set the data variables typesetting the letter
only takes a few lines of code.

\starttyping
\starttext
\setups [letter:place]
\stoptext
\stoptyping

Interfacing to the \XML\ database takes the following steps:

\starttyping
\setvariables
  [letter:data]
  [   address=\XMLflush{address},
       prefix=\XMLflush{prefix},
       suffix=\XMLflush{suffix},
     initials=\XMLflush{initials},
   formalname=\XMLflush{formalname},
 informalname=\XMLflush{informalname}]
      content=\getbuffer]
\stoptyping

Here we assume that the content goes into a buffer. Of
course we need to set the selection and filename for the
contacts handler:

\starttyping
\setvariables
  [contacts]
  [selection=...,
   file=...]
\stoptyping

Next we hook the letter style into the contacts handler:

\starttyping
\startsetups[contact:handle]
  \setups[letter:place]
\stopsetups
\stoptyping

And now we're ready to process the database and turn the
addresses that match our selection criterium into a
personalized letter.

\starttyping
\XMLprocesscontacts
\stoptyping

An example of this integration can be found in the \EXAMPLE\
module \type {x-e-0302.tex}.

\chapter{Examples}

On the following pages we will show a few examples of generating
letters.

\startbuffer[demo]

\startfiguretext
  {none}
  {\externalfigure
     [\getvariable{xcorresp}{demo}]
     [frame=on,
      width=.65\textwidth]}

  \switchtobodyfont[small]

  \typefile{\getvariable{xcorresp}{demo}}

\stopfiguretext

\stopbuffer

\page \setvariables[xcorresp][demo=letter-0] \getbuffer[demo]

In this example we choose a font and typeset a simple test
letter. We use a test setup to initialize the different
elements that make up a letter.

\page \setvariables[xcorresp][demo=letter-1] \getbuffer[demo]

One of the first things you will like to do is to add a logo
or other personal mark. There are several ways to do this.
Here we put a graphic in the background.

This option makes sense if your logo is placed on a paper
format similar to the one of your letter.

\page \setvariables[xcorresp][demo=letter-2] \getbuffer[demo]

There is a special layer \type {letterhead} that you can use
to position a graphic (or any content you want). You can
place multiple graphic (or text) elements.

\page \setvariables[xcorresp][demo=letter-3] \getbuffer[demo]

Layers permits you to finetune the position, for instance
with the \type {(v|h)offset} keys. The offsets are relative
to the corner. More information about layers can be found
in the \quote {details} manual.

\page \setvariables[xcorresp][demo=letter-4] \getbuffer[demo]

The adress and reference elements are also positioned in
layers, but this time they are not in the background but
flushed in the normal text stream. Their content tries to
honor the grid. Instead of manipulating the layers directly,
we use intermediate variables.

\page \setvariables[xcorresp][demo=letter-5] \getbuffer[demo]

The previous page showed some placement options but without
cues this is not that easy to do. Therefore we have turned
on some tracing options here.

\page \setvariables[xcorresp][demo=letter-9] \getbuffer[demo]

In addition to some style manipulation variables, we also
use variables for setting the actual content of the letter.

\page \setvariables[xcorresp][demo=letter-10] \getbuffer[demo]

This style is identical to the previous one, apart from one
setting:

\starttyping
\setvariables [letter:style] [alternative=ade]
\stoptyping

How are such alternatives implemented? When the module \type
{letter} (or file \type {m-letter.tex}) is loaded, this module
itself will start with loading \type {b-letter.tex} (before),
and end with loading \type {a-letter.tex} (after). Especially
the last one can overload and extend the style with local
preferences.

\page \setvariables[xcorresp][demo=letter-11] \getbuffer[demo]

In \type {a-letter.tex} there are sections marked as:

\starttyping
\startsetups[letter:alternative:ade]
  \setups[letter:ade:init,letter:common:style,letter:common:makeup]
\stopsetups
\stoptyping

These setups are expanded when a letter is generated. The next page
demonstrates that we have three such sections, which means that we
have three alternative choices.

\page \setvariables[xcorresp][demo=letter-12] \getbuffer[demo]

So, the magic is in defining a setup related to the alternative. The
last alternative (\type {cts}) is roughly defined as follows. The whole
patch of the basic style is encapsulated in a setup:

\starttyping
\startsetups[letter:alternative:cts]
  ....
\stopsetups
\stoptyping

Between these commands we first define the graphic. In this case, the
logo is taken from a company specific \METAPOST\ file:

\starttyping
\startreusableMPgraphic{common:graphic}
  loadfigure "mp-prag" number 501 xsized 4cm ;
\stopreusableMPgraphic
\stoptyping

We use the Computer Modern Mono Variable as main body font; this is
set up with:

\starttyping
  \usetypescript[modernvariable][\defaultencoding]
  \setupbodyfont[modernvariable,14.4pt]
\stoptyping

We also use hanging punctuation and are quite tolerant with respect
to paragraph building. We use whitespace between the paragraphs.

\starttyping
\setupalign
  [hanging]

\setuptolerance
  [verytolerant,stretch]

\setupwhitespace
  [big]
\stoptyping

Next we change the layout a bit:

\starttyping
\setuplayout
  [color=dark,
   width=middle,
   height=middle,
   topspace=.5cm,
   header=3cm,
   headerdistance=.5cm,
   bottomspace=1.5cm,
   footer=0pt,
   backspace=4.5cm,
   cutspace=1.5cm,
   leftmargin=3cm,
   margindistance=1cm]
\stoptyping

We start the first page pretty high:

\starttyping
\setupheader
  [state=high]
\stoptyping

We also (re)define a few colors:

\starttyping
\definecolor [dark] [s=0.3]
\definecolor [gray] [s=0.7]
\stoptyping

The logo (graphic) is put in its the main letter layer. Because we
want to repeat it, we change its state:

\starttyping
\setuplayer
  [lettermain]
  [state=repeat,
   hoffset=.5cm]

\setlayer
  [lettermain]
  [y=0cm,x=-.5cm]
  {\reuseMPgraphic{common:graphic}}

\setlayer
  [lettermain]
  [y=4cm]
  {\gray \scale
     [width=3cm,height=\dimexpr(\paperheight-7cm)]
     {\rotate{\tt CONTEXT TYPESETTING SERVICES}}}

\setlayer
  [lettermain]
  [preset=leftbottom,
   y=.5cm]
  {\dark \scale[width=3cm,height=2cm]{\tt\framed
     [offset=overlay,align=normal,frame=off]
     {\tt Ridderstraat 27-29\\
      8061 GH Hasselt NL\\
      +31(0)38 477 53 69\\
      www.pragma-cts.com}}}
\stoptyping

There ar etwo layers available: \type {lettermain} and \type
{letternext}. By default, the main layer is placed once, and
the next layer is places on following pages.

The following settings change some defaults of the letter style.

\starttyping
\setvariables
  [letter:reference]
  [line=1,
   noflines-min=2,
   noflines-max=5]

\setvariables
  [letter:address]
  [line=2,
   noflines=7]
\stoptyping

The sequence of placing content is as follows (in most cases there the
default sequence will do). Watch how we use local setups here.

\starttyping
\startlocalsetups [letter:sequence]

  \setups [letter:place:reference]  \endgraf
  \setups [letter:place:address]    \endgraf
  \setups [letter:place:opening]    \endgraf
  \setups [letter:place:content]    \page[no]
  \setups [letter:place:closing]    \endgraf
  \setups [letter:place:appendices] \endgraf

\stoplocalsetups
\stoptyping

\chapter{The GUI}

{\em to do}

\chapter{The Tools}

We are in the process of writing an associated tool for
sorting, cleaning up, and manipulating the correspondence
database. This tool will be part of the \CONTEXT\
distribution as soon as it's stable and documented.

\stoptext